\documentclass{article}
\usepackage{amsmath,amssymb}
\usepackage{graphicx}

\title{Martingale-Coupled Distributed Search on Bump Landscapes}
\author{}
\date{}

\begin{document}
\maketitle

\section{Introduction}

The consensus-coupled SGD framework of the companion paper proposes that
distributed learners should share discoveries through a soft consensus
field: each worker is weakly attracted toward the grand average of the
ensemble. This mechanism accelerates feature discovery by making one
worker's bump detection contagious to the rest. However, it has a
fundamental flaw: the consensus pull is always attractive. Every worker is
drawn toward the mean regardless of whether the mean contains any useful
structure. If enough workers happen to cluster near an arbitrary point in
the landscape, the attraction can amplify that cluster into a false
consensus. Workers coalesce not because they have discovered a bump, but
simply because they have attracted each other.

In this paper we propose a modification that eliminates this failure mode.
Instead of pulling every worker toward every other, we assign each pair a
random sign: half of the interactions are attractive and half are
repulsive. If worker $j$ is attracted to worker $i$, it receives a force
toward $j$'s position. If worker $j$ is repelled by worker $i$, it
receives a force away from $j$'s position. To prevent any deterministic
drift from accumulating, the signs are independently re-randomized at every single optimization step. The expected
force on any worker from the ensemble is therefore strictly zero at all times. This is the
martingale property.

The consequence is simple but powerful. Because the expected force is zero, the coupling no longer pulls agents toward each other. Instead, the coupling acts as a source of purely stochastic \emph{noise} (or effective temperature). In empty regions of the landscape where no bump is present, this injected noise is isotropic and bounded, allowing agents to diffuse freely and stably as if uncoupled. But when agents discover a bump and anchor themselves to it, the variance of the coupling force changes geometry. The anchored agents synthesize a massive, highly anisotropic thermal fluctuation that points exactly along the axis of the discovered bump. The exploring agents are not \emph{pulled} toward the bump; rather, the temperature in the direction of the bump becomes so hot that they are violently boiled over the regularizer's barrier. 

The system therefore has a built-in filter: it amplifies real structure
through directed temperature and ignores false consensus. This preserves the discovery-propagation mechanism while eliminating the pathological
attraction to empty space.

\section{Setup}

We use the bump data-generating process from the companion papers. The
regression function $f \colon \mathbb{R}^d \to \mathbb{R}$ is
\begin{equation}\label{eq:dgp}
  f(x) = \sum_{i=1}^{k} \frac{b_i}{1 + (x - c_i)^\top A_i\,(x - c_i)},
\end{equation}
where each bump has height $b_i$, center $c_i$, and positive
semi-definite shape matrix $A_i$.

We have $n$ agents with positions
$\theta_1, \dots, \theta_n \in \mathbb{R}^d$. At each time step $t$, we draw
a fresh sign matrix $S^t \in \{-1, 0, +1\}^{n \times n}$ with the following
properties:
\begin{itemize}
  \item $S_{ii}^t = 0$ (no self-interaction),
  \item For all $i \neq j$, $S_{ij}^t$ is drawn independently and uniformly from $\{-1, +1\}$.
\end{itemize}
Because the signs are fully independent (IID) and re-drawn every step, the expected coupling force on any exploring agent is exactly zero. There is no coherent drift.

In a low-dimensional setting, recruiting the entire swarm into a single feature would be catastrophic, as the landscape outside the bump would become a ``ghost town'' devoid of explorers. However, in our overparameterized regime where the ambient dimension is vast ($d \sim 10^9$ to $10^{12}$) and features are low-rank (e.g., $r = 2$), recruiting the entire ensemble to a discovered bump is actually desirable. The agents converge only in the $r$-dimensional subspace defining the bump. In the remaining $d - r$ orthogonal dimensions, the agents remain functionally independent and free to explore. The vast null space ensures that exploration continues unabated even if the entire swarm is recruited to exploit a specific low-rank feature.

The update rule for agent $i$ is
\begin{equation}\label{eq:martingale-update}
\theta_i^{t+1}
= \theta_i^t + \eta\,g_i(\theta_i^t)
+ \gamma \sum_{j \neq i} S_{ij}^t\,(\theta_j^t - \theta_i^t)
+ \xi_i^t,
\end{equation}
where $g_i$ is the local gradient (as ascent on $f$), $\gamma > 0$ is the
coupling strength, and $\xi_i^t \sim \mathcal{N}(0, 2D)$ is mini-batch noise.

\section{The Martingale Property and Null-Space Stability}

Consider a single agent $i$ exploring the null space, held near the origin by a global regularizer $\lambda \|\theta\|^2$. The coupling force on agent $i$ is
\[
F_i^t = \gamma \sum_{j \neq i} S_{ij}^t\,(\theta_j^t - \theta_i^t).
\]
Because $S_{ij}^t$ is re-randomized at every step, it is perfectly uncorrelated with the positions $\theta^t$. Therefore, $\mathbb{E}[F_i^t] = 0$. The force
is a strict martingale: it has zero drift. 

Because there is no drift, the martingale coupling acts purely as white noise. If the other agents are wandering near the origin with variance $\sigma^2$, the injected noise $F_i$ has variance proportional to $\gamma^2 n \sigma^2$. As long as $\gamma$ is kept small enough that this injected variance does not overwhelm the restoring force of the regularizer, the null space is unconditionally stable. The agents will never experience a runaway explosion or self-reinforcing false coalescence. They simply experience a slight increase in effective temperature in the null dimensions and continue to diffuse freely.

\section{Contrast with Consensus Coupling}

In the original consensus-coupled SGD, every agent is attracted to the
grand mean $\bar{\theta}$. The force on agent $i$ is
$-\gamma(\theta_i - \bar{\theta})$, which is always directed toward the
center of mass. This has two failure modes:

\paragraph{False coalescence.}
If a subset of agents happens to cluster in empty space (by stochastic
fluctuation), the resulting shift in $\bar{\theta}$ creates a persistent, deterministic drift that pulls all other
agents toward that cluster. The cluster grows not because it contains
structure but because the attraction is self-reinforcing. In the
martingale scheme, a random cluster produces mean-zero noise. Without deterministic drift, the random cluster rapidly dissolves.

\paragraph{Averaging trap.}
When two agents discover two different bumps, the grand mean lies between
them in empty space. Both agents are pulled away from their bumps toward
the void. In the martingale scheme, there is no central average
that pulls everyone to the middle. The forces remain mean-zero and simply increase local variance.

\section{Simulation}

We test the martingale mechanism on a minimal example: 100 agents in
$\mathbb{R}^2$ with a single bump. Agent~0 starts at the bump center;
the remaining 99 agents are initialized uniformly at random in
$[-50, 50]^2$. The bump is located at $(30, 30)$ with height $b = 5$
and shape $A = 0.5\,I_2$.

We compare two conditions:
\begin{enumerate}
  \item \textbf{Random search}: agents perform independent noisy gradient
    ascent with no inter-agent coupling ($\gamma = 0$).
  \item \textbf{Martingale coupling}: agents perform the same gradient
    ascent but with the signed coupling force from
    equation~\eqref{eq:martingale-update}, re-randomizing $S^t$ at every step.
\end{enumerate}
Discovery is defined as any agent besides agent~0 entering within
distance~$2$ of the bump center. We run 200 independent trials for each
condition and record the discovery time.

\begin{figure}[ht]
  \centering
  \includegraphics[width=\textwidth]{fig_martingale_discovery.pdf}
  \caption{Left: histogram of discovery times across 200 trials.
    Right: cumulative fraction of trials in which the bump has been
    discovered by step $t$. Martingale coupling substantially accelerates
    discovery compared to independent random search.}
  \label{fig:discovery}
\end{figure}

\begin{figure}[ht]
  \centering
  \includegraphics[width=\textwidth]{fig_martingale_snapshot.pdf}
  \caption{Spatial snapshot of agent positions at $t = 300$ for a single
    trial. Left: random search, agents diffuse independently. Right:
    martingale coupling, agents near the bump cluster while others
    continue exploring. The red star marks the discoverer (agent~0); the
    dashed circle marks the discovery radius.}
  \label{fig:snapshot}
\end{figure}

\section{The Convex Case: Swarm Variance and the Critical Threshold}

Before analyzing the escape dynamics in a multimodal landscape, it is instructive to consider the global behavior of the swarm in the convex null space, governed solely by the global regularizer $f(\theta) = \frac{1}{2} \theta^\top H \theta$ (where $H = 2\lambda I$). 

When $n$ agents run independent SGD on a quadratic, they form an Ornstein-Uhlenbeck (OU) process. The stationary covariance of a single uncoupled agent is $\Sigma_{\text{sgd}} \approx \frac{\eta}{2} H^{-1} C$, where $C$ is the gradient noise covariance for a mini-batch of size $B$ \cite{mandt2017stochastic}. For $n$ independent workers, the variance of the ensemble mean $\bar{\theta} = \frac{1}{n} \sum \theta_i$ drops linearly to $\frac{1}{n} \Sigma_{\text{sgd}}$, perfectly simulating a single large batch of size $nB$.

Introducing the martingale coupling fundamentally alters this steady-state distribution. If we strictly enforce Newton's Third Law in the sign matrix---drawing symmetric signs $S_{ij}^t = S_{ji}^t$ such that every attractive pull is matched by an equal and opposite pull---the sum of all coupling forces across the swarm is exactly zero ($\sum_{i=1}^n F_i^t = 0$). Consequently, the ensemble mean $\bar{\theta}$ is completely blind to the martingale coupling; it evolves exactly as it would under standard data-parallel SGD, retaining the perfect $1/n$ variance reduction.

However, the individual agents $\theta_i$ experience the coupling as a massive source of injected noise. As derived earlier, the variance of the force injected into agent $i$ by the rest of the swarm is proportional to $\gamma^2 (n-1) \Sigma_{\text{mc}}$, where $\Sigma_{\text{mc}}$ is the stationary covariance of the coupled agents. 

Because this martingale noise is independent of the mini-batch gradient noise, the total effective noise covariance driving agent $i$ is $C_{\text{total}} = \eta^2 C + \gamma^2 (n-1) \Sigma_{\text{mc}}$. Solving the Lyapunov equation $\eta H \Sigma_{\text{mc}} + \eta \Sigma_{\text{mc}} H = C_{\text{total}}$ yields the exact stationary covariance of the martingale swarm:
\begin{equation}
  \Sigma_{\text{mc}} = \left( 2\eta H - \gamma^2 (n-1) I \right)^{-1} \eta^2 C.
\end{equation}

This reveals a profound structural constraint on distributed search. In a purely convex landscape, the martingale coupling is a strict penalty: the individual agents are artificially heated, expanding their stationary variance ($\Sigma_{\text{mc}} > \Sigma_{\text{sgd}}$) in order to hunt for non-convex features. 

More importantly, the denominator dictates an absolute explosion threshold. For the covariance to remain positive definite and stable, we must satisfy:
\begin{equation}
  \gamma^2 (n-1) < 2\eta \lambda_{\min}(H).
\end{equation}
If the swarm size $n$ or the coupling strength $\gamma$ grows too large, the injected communication noise overwhelms the regularizer's restoring force, and the swarm geometrically explodes to infinity.

This mirrors the concept of the \emph{critical batch size} in distributed training \cite{mccandlish2018empirical, shallue2019measuring, jain2018parallelizing}. Just as there is a critical batch size beyond which adding more data-parallel workers yields diminishing returns because the gradient noise falls below the landscape curvature \cite{jain2018parallelizing}, the martingale swarm has a \textbf{critical swarm size} $n_{\text{crit}} \approx 2\eta \lambda_{\min} / \gamma^2$. If $n$ exceeds this threshold, the swarm can no longer be contained by the convex basin, and the search process collapses.

\section{Anisotropic Temperature and Directional Escape}

To understand why pure white noise accelerates discovery rather than just making the system chaotic, we must examine the covariance matrix of the injected noise.

Suppose we want to speed up exploration by simply increasing the learning rate or injecting more standard SGD noise (raising the effective temperature $D$). In a $d$-dimensional space where $d \sim 10^9$, isotropic noise spreads equally in all directions. The probability that an isotropic random walk happens to take a sequence of steps specifically along the low-dimensional subspace of a distant bump is vanishingly small. Isotropic noise heats up the vast, empty null space just as much as the relevant feature dimensions, worsening the curse of dimensionality.

The martingale coupling entirely avoids this through \emph{anisotropic subspace noise}. The coupling force $F_i^t$ is strictly a linear combination of the agents' current positions. Its behavior depends entirely on where the other agents are:
\begin{enumerate}
  \item \textbf{Exploration Phase (Wilderness Clutter):} If all other agents are independently wandering near the origin with small variance $\sigma^2$, the covariance of the injected noise is roughly $\gamma^2 n \sigma^2 I_{n-1}$. The injected noise is isotropic but extremely small, making it harmless in the null space.
  \item \textbf{Exploitation Phase (The Directional Hot Plate):} Suppose $k$ agents discover a bump and become anchored at its center $c$, far from the origin ($\|c\| \gg \sigma$). The coupling force on agent $i$ now contains a massive term $\gamma (\sum_{j \in B} S_{ij}^t) c$. Because the signs are IID coin flips, this sum is a zero-mean scalar multiplier with variance $k$.
\end{enumerate}

The resulting covariance matrix of the noise injected into the exploring agents develops a massive rank-1 spike:
\begin{equation}
  \Sigma_{\text{injected}} = \mathbb{E}[F_i^t (F_i^t)^\top] \approx \gamma^2 k \, (c c^\top) + \text{clutter}.
\end{equation}
The swarm has effectively performed decentralized Principal Component Analysis. By locking $k$ agents into the bump, the ensemble synthesizes a highly anisotropic noise profile that points \emph{exactly} along the axis connecting the origin to the discovered feature. 

From the perspective of Freidlin-Wentzell large deviations theory, the expected time to escape the global regularizer's basin and reach the bump scales exponentially with the barrier height divided by the temperature: $\exp(\Delta U / D_{\text{eff}})$. 
In the orthogonal null dimensions, the temperature remains cold ($D_{\text{eff}} \approx D$). But along the specific 1D axis leading to the bump, the effective temperature becomes massively hot: $D_{\text{eff}} \approx D + \frac{1}{2} \gamma^2 k \|c\|^2$.

The martingale coupling does not lower the barrier. Instead, it creates a directional hot plate. The exploring agents are violently boiled over the regularizer's wall towards the bump, while the vast $10^{12}$ orthogonal dimensions remain cold and stable.

\subsection{The Thermodynamic Race ($\tau = 1$ limit)}

While injecting pure white noise ($\tau = 1$) provides an exponential speedup for exploring agents, it simultaneously exposes the anchored agents to massive thermal fluctuations. Because the anchored agent is pulled by $n-1$ exploring agents located near the origin, the variance of the force it experiences is $n-1$ times larger than the variance it injects into any single explorer. 

This creates a thermodynamic race. For the swarm to succeed, the exploring agents must be boiled out of the regularizer's basin \emph{before} the anchored agent is violently knocked out of the bump.

Let $E_{\text{reg}} = \lambda \|c\|^2$ be the energy barrier of the regularizer. The expected time for an explorer to reach the bump is:
\begin{equation}
  \tau_{\text{explore}} \propto \exp\left( \frac{E_{\text{reg}}}{\frac{1}{2} \gamma^2 \|c\|^2} \right).
\end{equation}
Conversely, if the bump has depth $b$, the expected time for the anchored agent to be knocked out by the $n-1$ explorers is:
\begin{equation}
  \tau_{\text{anchor}} \propto \exp\left( \frac{b}{\frac{1}{2} \gamma^2 (n-1) \|c\|^2} \right).
\end{equation}

For the cascade to grow, we require $\tau_{\text{explore}} \ll \tau_{\text{anchor}}$. Because the coupling parameter $\gamma^2$ cancels identically from both exponents, the viability of the pure-noise ($\tau=1$) martingale search simplifies to a strict, parameter-free structural condition:
\begin{equation}
  b > (n-1) E_{\text{reg}}.
\end{equation}

If a feature is sufficiently deep (its local depth exceeds $n-1$ times the global regularization penalty at that location), the anchored agent is unconditionally stable against the swarm's noise. The martingale coupling can then be safely tuned to maximize the exponential speedup $\tau_{\text{explore}}$, ensuring rapid discovery of the structure without risking runaway explosion in the vast null space.

\section{Multiple Bumps and Exploration--Exploitation}

When the landscape contains multiple bumps, the martingale mechanism
naturally partitions the variance. The agents anchored in bump~1 synthesize noise pointing along the axis to bump~1, while agents in bump~2 synthesize noise pointing to bump~2. The total injected covariance matrix becomes a low-rank superposition of the discovered feature directions.

This provides automatic load balancing and an elegant exploration-exploitation tradeoff without any central coordination. The swarm dynamically detects the axes of relevant features and selectively turns up the temperature along those specific axes, accelerating the recruitment of exploring agents into known structure while preserving a cold, stable null space for continued search.

\section{Open Questions}

\begin{itemize}
  \item How does the nucleation cascade rate depend on $n$, $d$, and the
    bump geometry?
  \item Should the sign matrix $S$ be symmetric ($S_{ij}^t = S_{ji}^t$) or
    asymmetric? Symmetric signs create Newton's-third-law pairs;
    asymmetric signs allow richer interaction patterns.
  \item Can the sign matrix be adapted online based on agent fitness, or
    does IID assignment suffice?
  \item How does the martingale mechanism interact with the filtered
    center $c_t$ from the consensus paper?
  \item In realistic neural network training, what is the analogue of the
    sign matrix? Can it be implemented efficiently at scale?
\end{itemize}

\bibliographystyle{plain}
\bibliography{distributed_bumps}

\end{document}